\chapter{System Design}
\label{chap:design}
%=====================================================================
In the diagrams of this chapter, solid elements are implemented on the
\texttt{develop} branch and covered by tests, elements tagged ``partial'' are
implemented slices, dashed elements are planned, and red marks a known security
gap. All diagrams were drawn from the code at commit \texttt{ce49dcc}
\tbd{(update to the tagged submission commit)}; each figure's source file in
\path{ThesisFigs/src/} names the code it was checked against.

\section{System Architecture}
\label{sec:architecture}
Aegis is organised in five layers (Figure~\ref{fig:layers}). Layers talk only to
their neighbours, through files on the Modal volume or through HTTPS JSON.

\aegisfigtall{ch5-layered-architecture}
  {Layered architecture of Aegis at the end of FYP-1. Colour marks the layer;
   dashed boxes are planned; red marks a security gap.}
  {fig:layers}

\textbf{Capture and ingestion} runs on the laptop. \texttt{aegis-ingest} reads
the phone's MJPEG stream, samples it at 2\,fps, keeps only sharp, well-exposed
frames that differ from the last kept frame, and writes JPEG frames with an
index of frame and gap records. A finished or growing capture is cut into
segments of at most 120 frames and uploaded to a Modal volume. Batch videos
follow the same path through \texttt{ingest\_video}.

\textbf{3D processing and perception} runs on Modal. A CPU worker normalises
the frames; an L4 GPU worker runs VGGT-$\Omega$ to recover camera poses and
per-pixel 3D points, and a RANSAC floor fit finds ``up''; a second GPU worker
detects objects with SAM~3 (or YOLOE), refines masks with SAM~2, tracks
identities with DINOv3 embeddings, and lifts each detection into a
gravity-aligned oriented box. With \texttt{AEGIS\_GPU=cpu} the same stages run
on CPU containers.

\textbf{Memory} turns boxes into rooms, anchors and objects, classifies each
object's relation to its anchor, debounces observations into state runs and
derives events, producing one session graph. A session can then be merged into
a per-home SQLite store that keeps stable identities, groups activities, keeps
evidence crops and compacts history older than three days.

\textbf{Query service} is a FastAPI application deployed on Modal. It answers
ranked object queries (\path{/query}) and natural-language questions
(\path{/ask}) from the session graph or the home store. It currently allows
any origin and has no authentication on writes.

\textbf{Presentation} is the Patient and Caregiver apps, the \texttt{aegis-ask}
terminal client, and an offline Rerun viewer.

\section{Architectural Components}
\label{sec:components}
Figure~\ref{fig:components} shows one component per Python package plus the
three clients, grouped by where they run. Dashed arrows are package imports
(«use»), taken from the code. The only network interface is the HTTP API,
drawn as a ball (provided by \texttt{remote.api}) and a socket (required by
each client). Table~\ref{tab:components} lists what each component provides and
requires.

\aegisfig{ch5-components}
  {UML component diagram: Python packages, Modal apps and clients. Dashed
   arrows are imports; the ball and socket is the HTTP API.}
  {fig:components}

\begin{table}[H]\centering\footnotesize
\caption{Component responsibilities.}\label{tab:components}
\begin{tabularx}{\textwidth}{@{}L{2.0cm} L{3.4cm} Y L{2.8cm}@{}}\toprule
\rowcolor{hdrfill}\textbf{Component} & \textbf{Package path} & \textbf{Provides} & \textbf{Requires}\\\midrule
core & \path{src/aegis/core/} & \texttt{RuntimeConfig}, \texttt{Vocabulary}, \texttt{DatasetPaths}, models & --\\
ingestion & \path{src/aegis/ingestion/} & \texttt{FrameSource}, \texttt{ingest\_live}, \texttt{promote\_live\_segments} & core\\
geometry & \path{src/aegis/geometry/} & \texttt{run\_vggt\_reconstruction}, \texttt{GravityFrame}, \texttt{lift\_detection\_to\_3d} & core\\
perception & \path{src/aegis/perception/} & \texttt{detect\_and\_track\_objects} & core\\
memory & \path{src/aegis/memory/} & \texttt{MemoryRepository}, \texttt{HomeStore}, \texttt{ingest\_session}, \texttt{compact} & core, geometry\\
nlq & \path{src/aegis/nlq/} & \texttt{AnswerEngine.ask}, \texttt{DictionaryIndex} & memory (graph view)\\
remote.workers & \path{remote/workers.py}, \texttt{modal\_app.py} & \texttt{run\_preprocess}, \texttt{run\_vggt}, \texttt{run\_detect\_and\_lift}, \texttt{run\_rerun} & ingestion, geometry, perception, memory, visualization\\
remote.api & \path{remote/api.py}, \texttt{serve.py} & HTTP API \path{/api/v1} & memory, nlq\\
visualization & \path{visualization/rerun_scene.py} & \texttt{run\_rerun} & graph and geometry artifacts\\
evaluation & \path{src/aegis/evaluation/} & readiness and score reports & memory, artifacts\\
Patient App & \texttt{apps/Patient App/} & -- & HTTP API\\
Caregiver App & \texttt{apps/Caregiver App/} & -- & HTTP API\\
\texttt{aegis-ask} & \path{src/aegis/nlq/cli.py} & -- & HTTP API or local graph\\
\bottomrule\end{tabularx}\end{table}

The prototype is deployed across six kinds of node (Figure~\ref{fig:deployment}).
The wearer's phone only streams video. The laptop runs capture, launches the
pipeline and, when artifacts are downloaded, holds the persistent home store.
Modal runs CPU and GPU containers for processing and a separate CPU web
endpoint, \texttt{aegis-memory-api}, that serves the API from the data volume,
so normal operation needs no download to the laptop. The patient's and
caregiver's phones reach that endpoint over HTTPS. GitHub Actions runs linting,
the test suite with a 75\% coverage floor and Flutter checks on every push to
\texttt{main} and \texttt{develop}, and publishes the patient APK. Model
weights are downloaded once from Hugging Face and cached on a model volume.

\aegisfigtall{ch5-deployment}
  {UML deployment diagram: devices, Modal execution environments, volumes and
   communication paths. Red: security gaps; dashed: planned or development only.}
  {fig:deployment}

\section{Interfaces and Dependencies}
\label{sec:interfaces}
Table~\ref{tab:api} is the HTTP API contract (\path{src/aegis/remote/api.py},
base path \path{/api/v1}, \texttt{API\_VERSION = 0.2.1}). There are 26 route
handlers; six are simulator presets.

\begin{table}[H]\centering\scriptsize
\setlength{\tabcolsep}{3pt}
\caption{HTTP API contract.}\label{tab:api}
\begin{tabularx}{\textwidth}{@{}L{0.9cm} L{3.2cm} L{2.2cm} Y L{2.0cm} L{1.6cm}@{}}\toprule
\rowcolor{hdrfill}\textbf{Method} & \textbf{Path} & \textbf{Inputs} & \textbf{Response} & \textbf{Mutates} & \textbf{Auth}\\\midrule
GET & \path{/health} & \texttt{dataset\_id?}, \texttt{home\_id?} & \texttt{HealthResponse} (anchors, rooms, events, frames, \texttt{compaction\_ok}) & no & none\\
GET & \path{/scene_graph} & \texttt{dataset\_id?}, \texttt{home\_id?} & full graph JSON & no & none\\
GET & \path{/datasets} & -- & datasets and tracked objects & no & none\\
POST & \path{/dataset/select} & \texttt{dataset\_id} & active dataset & \textbf{yes} & none\\
GET & \path{/anchors} & \texttt{dataset\_id?}, \texttt{home\_id?} & \texttt{AnchorModel[]} & no & none\\
GET & \path{/rooms} & \texttt{dataset\_id?}, \texttt{home\_id?} & \texttt{RoomModel[]} & no & none\\
POST & \path{/rooms/rename} & \texttt{room\_id}, \texttt{label} & renamed room & \textbf{yes} (graph, store, \texttt{home\_map.json}) & none\\
POST & \path{/enroll} & \texttt{EnrollRequest} & personal items & \textbf{yes} (store) & none\\
GET, POST & \path{/ask} & \texttt{text} ($\leq$ 300 chars) & \texttt{AskResponse} & no & none\\
GET & \path{/query} & \texttt{item} & \texttt{MemoryQueryResponse} ($\leq$ 3 candidates) & no & none\\
GET & \path{/timeline} & \texttt{start?}, \texttt{end?} & \texttt{TimelineEventModel[]} (with \texttt{tier}) & no & none\\
GET & \path{/activities} & \texttt{home\_id}, \texttt{start?}, \texttt{end?} & \texttt{ActivityModel[]} & no & none\\
GET & \texttt{/crops/\{instance\_id\}} & -- & JPEG & no & none\\
GET & \texttt{/video/\{dataset\_id\}} & -- & MP4 & no & \textcolor{red}{none; raw video}\\
GET & \path{/capture/status} & -- & \texttt{CaptureStatusModel} & no & none\\
POST & \path{/capture/pause}, \path{/resume} & -- & \texttt{CaptureStatusModel} & \textbf{yes} (pause file) & none\\
POST, GET & \path{/simulate_transition}, \path{/reset}, \path{/simulate/*} (6) & -- & simulator responses & \textbf{yes} & gated, off by default\\
\bottomrule\end{tabularx}\end{table}

Table~\ref{tab:artifacts} lists the file contracts between layers
(\path{docs/artifact_contracts.md}). All timestamps are UTC ISO strings
(\path{core/timeutil.py}).

\begin{table}[H]\centering\footnotesize
\caption{Artifact contracts.}\label{tab:artifacts}
\begin{tabularx}{\textwidth}{@{}L{4.2cm} l Y@{}}\toprule
\rowcolor{hdrfill}\textbf{Artifact} & \textbf{Schema} & \textbf{Producer $\rightarrow$ consumer}\\\midrule
\texttt{captures/<id>/index.jsonl} (frame and gap records) & 2 & \texttt{ingest\_live} $\rightarrow$ promotion\\
\texttt{session.json}, \texttt{session\_state.json} & 2 & live writer $\rightarrow$ status, crash resume\\
\texttt{metadata.json} & 2 & preprocess / promotion $\rightarrow$ reconstruction\\
\texttt{poses.json}, \path{pointmaps/*.npz}, \texttt{gravity.json} & -- & \texttt{run\_vggt} $\rightarrow$ perception\\
\texttt{detected\_objects\_3d.json} & 2 & perception $\rightarrow$ memory build\\
\texttt{r4d\_scene\_graph.json}, \texttt{r4d\_timeline.json} & 2 & memory build $\rightarrow$ API, home store, Rerun\\
\texttt{events.json}, \texttt{spatial\_memory\_graph.json} & -- & memory build $\rightarrow$ evaluation\\
\texttt{memory.db} & store 1 & \texttt{ingest\_session} $\rightarrow$ API (\texttt{home\_id})\\
\texttt{home\_map.json} & -- & room rename $\rightarrow$ room matching\\
\bottomrule\end{tabularx}\end{table}

\section{Architecture Decisions and Rationale}
\label{sec:adr}
Table~\ref{tab:adr} records the main architecture decisions, the alternatives
considered, why each was chosen and the evidence.

\begin{table}[H]\centering\scriptsize
\setlength{\tabcolsep}{3pt}
\caption{Architecture decision record.}\label{tab:adr}
\begin{tabularx}{\textwidth}{@{}l L{3.0cm} L{2.6cm} Y L{2.6cm}@{}}\toprule
\rowcolor{hdrfill}\textbf{\#} & \textbf{Decision} & \textbf{Alternatives} & \textbf{Rationale} & \textbf{Evidence}\\\midrule
AD1 & Raw VGGT-$\Omega$ coordinates; gravity stored beside them & Rotate or scale the scene; the mid-era SCALE/\texttt{to\_zup} transform (removed) & No hidden transform; every derived height is reproducible & \path{geometry/gravity.py}; \texttt{test\_gravity} (12)\\
AD2 & Rooms from wearer trajectory and sparse landmarks & Dense point-cloud segmentation & Cheap; works on RGB-only video & \path{memory/rooms.py}; Hydra, HOV-SG\\
AD3 & Debounced state runs instead of a per-frame timeline & Store every observation & Graph size independent of video length; 5.7--7.2$\times$ smaller & \path{memory/runs.py}; Section~\ref{sec:compaction-audit}\\
AD4 & Per-home SQLite store & Growing JSON; a database server & Single file, transactional, no service to run & \path{memory/store.py}\\
AD5 & Three-day detail, verified template summaries, roll-up & Keep everything; LLM summaries & Bounded storage; no model can hallucinate a summary; failure keeps detail & \path{memory/compaction.py}\\
AD6 & SAM~3 default detector & YOLOE, YOLO-World & Only model that found small personal objects and hands & Section~\ref{sec:experiments}\\
AD7 & DINOv3 re-identification & HSV histograms & Appearance survives viewpoint and lighting change & \path{perception/detector.py}\\
AD8 & Closed-vocabulary rule engine; LLM only as an optional parser & A VLM or LLM writes answers & Never states an unobserved location & \path{nlq/engine.py}\\
AD9 & Artifacts stay on the Modal volume; API served there & Download to the laptop and serve locally & No large transfers; phones reach HTTPS & \path{remote/serve.py}\\
AD10 & Thin apps that never mutate memory (except caregiver rename and enrol) & On-device processing & Simple, battery-friendly; one source of truth & \texttt{CONTEXT.md} invariant 7\\
AD11 & Ranked top-3 candidates with photos & A single answer & Honest about ambiguity & \texttt{reasoning.rank\_candidates}\\
\bottomrule\end{tabularx}\end{table}

\section{Design Evolution}
\label{sec:evolution}
Figure~\ref{fig:evolution} compares the design at the mid report (commit
\texttt{2298112}, 17~October) with the FYP-1 final design. Each changed element
is numbered, and Table~\ref{tab:evolution} gives the reason for each change.

\aegisfigtall{ch5-design-evolution}
  {Design evolution from the mid report to the FYP-1 final design. Orange
   borders mark changed elements; the numbers match Table~\ref{tab:evolution}.}
  {fig:evolution}

\begin{table}[H]\centering\footnotesize
\caption{Design changes since the mid report and their reasons.}\label{tab:evolution}
\begin{tabularx}{\textwidth}{@{}l Y Y@{}}\toprule
\rowcolor{hdrfill}\textbf{\#} & \textbf{Change} & \textbf{Why}\\\midrule
1 & YOLOE / YOLO-World $\rightarrow$ SAM~3 (demo) / YOLOE (live) + SAM~2 + DINOv3 re-ID & Detector benchmark (Section~\ref{sec:experiments})\\
2 & Raw VGGT coordinates $\rightarrow$ + gravity frame & Relations need ``up''\\
3 & Per-frame timeline (schema 1) $\rightarrow$ state runs and latest states (schema 2) & Storage growth\\
4 & Three relations $\rightarrow$ nine relations in the gravity frame & ``Inside / under'' questions\\
5 & JSON graph only $\rightarrow$ + SQLite home store and cross-session matching & Multi-session memory and restarts\\
6 & Compaction planned $\rightarrow$ compaction and roll-up implemented & Proposal action 1\\
7 & \texttt{QueryReasoningEngine} only $\rightarrow$ + \texttt{rank\_candidates} and the \path{/ask} NLQ engine & Ambiguity and natural language\\
8 & Local FastAPI $\rightarrow$ Modal-deployed \texttt{aegis-memory-api} & Phone access without a laptop\\
9 & Single \path{apps/mobile} $\rightarrow$ Patient App + Caregiver App; live sampler joined to the pipeline & Panel action 4; the live-to-memory gap closed\\
\bottomrule\end{tabularx}\end{table}

Several elements were also \emph{removed}: the SCALE=5.0/\texttt{to\_zup}
display transform of the mid prototype; the Depth-Anything depth model and dead
DROID code (\texttt{c4236df}); the client-side mock answer (\texttt{180b897});
the legacy \path{apps/mobile} app; and the separate \texttt{NL Querying/}
folder, now archived after its consolidation into \path{src/aegis/nlq}.

\section{Domain Model}
\label{sec:domain}
The domain model (Figure~\ref{fig:domain}) separates what persists from what is
derived. A \textbf{Home} owns \textbf{Rooms} and \textbf{Sessions}.
Furniture-sized \textbf{Anchors} sit in rooms. Each personal
\textbf{ObjectInstance} has a history of \textbf{StateRuns}, each a stable
relation to an anchor over a time interval, and \textbf{Events} record the
change between two runs. \textbf{Activities} group the wearer's stay in a room
with the events that happened there; \textbf{RoomVisits} record each stay.
\textbf{Evidence} is a photo crop of a run, and \textbf{Summaries} replace
detail older than three days. Rooms, anchors, sessions and each object's
last-known state are never compacted, so ``where is X now?'' survives any
amount of history. Caregiver access (caregiver, access grant and audit entry)
is planned and drawn dashed. The enumerated relation and event types are
listed in Section~\ref{sec:relations-algo}.

\aegisfig{ch5-domain-model}
  {Conceptual domain model (attributes only). The dashed box is planned
   caregiver access (FR8).}
  {fig:domain}

\section{Data Design}
\label{sec:data}
\textbf{Session graph.} Each processed session produces one
\texttt{r4d\_scene\_graph.json} (schema version 2). Figure~\ref{fig:scenegraph}
shows its top-level keys. The file holds no per-keyframe data: an object
standing still for an hour is one state run, so the file does not grow with
video length. Arrows are id references inside the document, not database
foreign keys.

\aegisfig{ch5-scene-graph-v2}
  {Session scene graph (schema 2): one box per top-level key; arrows are id
   references.}
  {fig:scenegraph}

\textbf{Home store.} The persistent memory of one home is a single SQLite file,
\texttt{<store\_root>/<home\_id>/memory.db} (\texttt{SCHEMA\_VERSION = 1}),
with twelve tables. Figure~\ref{fig:erd-core} shows the core entities and
Figure~\ref{fig:erd-events} the event, visit, activity and retention tables.
Every column is copied from \path{memory/store.py}. Primary keys are bold and
foreign keys are italic; some id columns are soft references with no
foreign-key constraint, marked ``(soft)'', so that compaction can delete
detail rows without cascading. Colour shows the retention class: green tables
(\texttt{room}, \texttt{session}, \texttt{anchor}, \texttt{instance},
\texttt{meta}) are never compacted; orange tables (\texttt{state\_run} except
protected last-known runs, \texttt{event}, \texttt{room\_visit},
\texttt{activity}, \texttt{evidence}) are compacted after three days; blue
tables (\texttt{summary}, \texttt{compaction\_log}) are written by compaction.
\texttt{meta} holds the schema version and the personal-items JSON.

\aegisfigtall{ch5-erd-core}
  {Crow's-foot ERD of the home store, core entities.}
  {fig:erd-core}

\aegisfigtall{ch5-erd-events}
  {Crow's-foot ERD of the home store, events, visits, activities and
   retention tables. White boxes are shown in full in Figure~\ref{fig:erd-core}.}
  {fig:erd-events}

\textbf{Live capture contracts.} Figure~\ref{fig:artifacts} shows the files a
live capture writes and how promotion turns them into datasets. Each line of
\texttt{index.jsonl} is a frame record or a gap record; every field shown is
validated by \path{ingestion/contracts.py}. Promotion cuts the capture into
segments of at most 120 frames, splitting at long gaps, and each segment is
uploaded to the Modal volume.

\aegisfig{ch5-live-session-artifacts}
  {Live capture folder, the promoted segment dataset and the Modal volume.}
  {fig:artifacts}

\section{User Interface Design}
\label{sec:ui}
The panel asked for a product-like design suited to older users (proposal
action 4). The Patient App (Figure~\ref{fig:patient-flow}) opens on large
object tiles built only from objects that memory actually holds, a single large
microphone button and a call button. Answers appear in a bottom sheet with one
plain sentence, a labelled certainty badge (``Found it'', ``Best guess'', ``I'm
not sure --- which one?'', ``I'm not sure'') and up to three real photos, never
a bare probability. Both apps use 44-point touch targets and a warm
espresso/ochre theme.

\aegisfigtall{ch5-patient-app-flow}
  {Patient App screen flow. Labels are copied from the Flutter code;
   transitions name the API call they make.}
  {fig:patient-flow}

The Caregiver App (Figure~\ref{fig:caregiver-flow}) adds search history, a
tracking timeline, room naming, server and capture status, and personal-item
enrolment. Room rename and enrolment are writes without authentication; this is
shown in red and is planned work (FR8.1).

\aegisfigtall{ch5-caregiver-app-flow}
  {Caregiver App screen flow, with the four tabs, search history and the enrol
   dialog.}
  {fig:caregiver-flow}

\section{Process and Behaviour Design}
\label{sec:behaviour}
\textbf{Data flow.} Figure~\ref{fig:dfd0} shows Aegis as a single process with
its five external entities; there are no data stores at level~0.
Figure~\ref{fig:dfd1} decomposes it into seven processes and six data stores
(Gane--Sarson notation). Capture (1.0) writes the live folder (D1); promotion
(2.0) uploads segments to the Modal volume (D2); reconstruction and perception
(3.0) and the session-memory build (4.0) read and write D2, and 4.0 also writes
evidence crops (D5). Ingest and compaction (5.0) turns a downloaded graph into
home-store rows (D3) and the room registry (D4). Answering (6.0) reads D2, D3
and D5; management (7.0) writes caregiver room names, personal items and the
pause flag.

\aegisfig{ch5-dfd-level0}{Data flow diagram, level~0.}{fig:dfd0}
\aegisfigtall{ch5-dfd-level1}{Data flow diagram, level~1.}{fig:dfd1}

\textbf{Pipeline run.} Figure~\ref{fig:seq-pipeline} shows one run of
\texttt{run\_pipeline.py} on a live capture. For each segment the script
promotes and uploads the frames, then calls three Modal functions in turn; each
writes its outputs back to the volume. Only with \texttt{--download} are the
outputs fetched and ingested into the home store. Normally the artifacts stay on
the volume and the API serves them there.

\aegisfigtall{ch5-seq-pipeline}{Sequence diagram: pipeline run on Modal.}{fig:seq-pipeline}

\textbf{Patient query.} Figure~\ref{fig:seq-query} follows one question. Speech
is transcribed on the phone, so no audio leaves it. The app calls
\path{/query} first; the server ranks candidates from the latest object
states. If nothing is found it falls back to \path{/ask}, where the NLQ engine
answers or abstains with a reason. If no reply arrives within 8\,s the app shows
``Can't reach Aegis memory''; it never shows a guessed answer.

\aegisfigtall{ch5-seq-query}{Sequence diagram: patient asks where an object is.}{fig:seq-query}

\textbf{Caregiver operations.} Figure~\ref{fig:seq-caregiver} shows room
renaming, personal-item enrolment and capture pause and resume. All of these
writes are currently unauthenticated.

\aegisfigtall{ch5-seq-caregiver}{Sequence diagram: caregiver operations.}{fig:seq-caregiver}

\textbf{Home ingest.} Figure~\ref{fig:seq-ingest} shows how a session graph is
merged into the home store. An unchanged graph (same SHA-256) is skipped. A new
graph is matched room by room, then anchor by anchor, then object by object
(Section~\ref{sec:matching-algo}); state runs are inserted or extended inside
one transaction, events are written, and the room registry is exported. Old
days are then compacted.

\aegisfigtall{ch5-seq-home-ingest}{Sequence diagram: home-store ingest and compaction.}{fig:seq-ingest}

\textbf{Live capture session.} Figure~\ref{fig:state-live} is the state machine
of a capture. Sampling keeps frames that pass the keyframe test; a pause,
a storage stall and a lost stream each lead to their own state and back, and
every return records a gap. A session ends at the time limit or after 30\,s
without frames. Pressing Ctrl+C stops it from any state, and
\texttt{session\_state.json} is saved after every frame so a crashed capture can
resume.

\aegisfig{ch5-state-live-session}{State machine: live capture session.}{fig:state-live}

\textbf{Object instance.} Figure~\ref{fig:state-object} shows the states of one
object. A new state starts only after two consecutive frames agree (or one
final frame with confidence $\geq$ 0.5). Pickup and placement move between
\emph{Placed} and \emph{Carried}. Leaving view emits no event; the last-known
state is kept (object permanence). In the next session the object is either
found in the same place, extending its run, or elsewhere, which records
\texttt{moved\_while\_unobserved}.

\aegisfig{ch5-state-object}{State machine: object instance.}{fig:state-object}

\textbf{Compaction.} Figure~\ref{fig:compaction} shows the daily compaction
loop. Each local day older than three days is summarised by templates, and the
summary is verified against its source rows before any row is deleted. If
verification fails, the day keeps all its detail and the failure is logged.
Protected last-known runs are never deleted.

\aegisfig{ch5-activity-compaction}{Activity diagram: daily compaction with verification.}{fig:compaction}

\textbf{NLQ answer engine.} Figure~\ref{fig:nlq} shows how \path{/ask} decides.
Each check that fails returns a specific abstention reason code; only a
question about a known, observed item reaches an answer, which is composed from
a template and the most recent state (an object in view first).

\aegisfigtall{ch5-activity-nlq}{Activity diagram: NLQ answer engine.}{fig:nlq}
